\documentclass[10pt,a4paper]{article}

\usepackage[T1]{fontenc}
\usepackage[utf8]{inputenc}
\usepackage{lmodern}
\usepackage{amsmath,amssymb,amsthm}
\usepackage{booktabs,array}
\usepackage{geometry}
\usepackage{xcolor}
\usepackage{enumitem}
\usepackage{microtype}
\usepackage{hyperref}
\usepackage{fancyhdr}

\geometry{margin=1in}
\hypersetup{
  colorlinks=true,
  linkcolor=blue!55!black,
  citecolor=blue!55!black,
  urlcolor=blue!55!black,
  pdftitle={TOT-ally Forged: A Zero-Query Raw-Vector Forgery against TOT},
  pdfauthor={Martin Feussner}
}
\setlist{nosep,leftmargin=*}
\setlength{\parindent}{1em}
\setlength{\parskip}{0.25em}
\pagestyle{fancy}
\fancyhf{}
\fancyhead[L]{TOT-ally Forged}
\fancyhead[R]{13 October 2026}
\fancyfoot[C]{\thepage}
\renewcommand{\headrulewidth}{0.3pt}

\newtheorem{theorem}{Theorem}
\newcommand{\F}{\mathbb F}
\newcommand{\Verify}{\mathsf{Verify}}
\newcommand{\Sign}{\mathsf{Sign}}
\newcommand{\KeyGen}{\mathsf{KeyGen}}

\title{\bfseries TOT-ally Forged:\\A Zero-Query Raw-Vector Forgery against TOT}
\author{Martin Feussner\\
\small Selmer Center, University of Bergen\\
\small \texttt{martin.feussner@uib.no}}
\date{13 October 2026}

\begin{document}
\maketitle

\noindent\fcolorbox{black!50}{black!4}{\parbox{0.94\textwidth}{
\textbf{AI disclosure and review status.}
This audit used the OpenAI model Daybreak Blue at the ultra reasoning level.
OpenAI Codex performed AI-assisted source review, attack-history searching,
scheme reconstruction, cryptanalysis, implementation, experiments, and
technical writing.  Separate specialist AI agents performed adversarial
review and a fresh-context AI agent independently reproduced the result with
a separately written implementation.  These were AI review and AI
reproduction, not human verification.  The technical, practical,
adversarial AI, independent-reproduction, final attack-history, and
Publication Reviewer gates all passed before public release.  No independent
human verification was performed or required.  OpenAI Codex generated and
selected the title and intentionally chose ``TOT-ally Forged'' as wordplay
on the scheme name TOT.}}

\begin{abstract}
We give a public-key-only, zero-signing-query existential forgery against the
signature construction in Shen and Tang's \emph{TOT, a Fast Multivariate
Public Key Cryptosystem with Basic Secure Trapdoor}, IACR ePrint 2013/771.
The published signature interface takes a raw vector $M\in\F_2^{128}$ and
accepts a vector $\sigma$ precisely when the public quadratic map satisfies
$P(\sigma)=M$.  An adversary chooses any nonzero $\sigma$, evaluates the
public key, sets $M=P(\sigma)$, and outputs $(M,\sigma)$.  The ordinary
verifier accepts by construction, while the zero-query transcript makes $M$
fresh.  Because $M$ lies in the image of $P$, this is also an honestly
signable raw vector rather than an out-of-domain artifact.

We reconstruct and execute the attack at the same-inventors recommended
signature row $(q,n,d,t)=(2,128,6,64)$, for which more than $2^{90}$ security
is claimed.  The primary run generates an unselected conforming key,
serializes only its full 1,073,280-bit quadratic public map, and separates
key generation, the public-only forger, and the ordinary verifier.  The
forgery is accepted and four negative controls reject.  A separately written
fresh-context implementation independently reproduces the result.  The
claim is confined to the literal raw-vector message space specified by the
paper and same-inventors patent.  We make no claim against a separately added
document-hashing or preimage-resistant encoding layer, and no novelty claim
for signature-first or no-message existential forgery as a general technique.
\end{abstract}

\section{Introduction}

TOT is a multivariate public-key construction built from an extended-field
trapdoor and secret affine maps~\cite{tot}.  Its 2013 paper gives encryption
and signature algorithms.  For signatures, the public key is a system of
$n$ quadratic polynomials over $\F_q$, a message is an $n$-component vector,
and verification evaluates the public system on the alleged signature and
compares the result directly with that vector.

This interface permits the classical verification-first construction.  A
forger need not invert the public map: it may select a preimage first and use
public evaluation to learn its image.  With no prior signing query, the image
is a fresh message in the scheme's stated domain.  This is the familiar
signature-first or no-message weakness of bare trapdoor-inversion signatures; the
contribution here is its concrete application, source binding, and
full-parameter reproduction for TOT.

The result has the following properties:
\begin{itemize}
  \item it uses the ordinary equation printed for TOT verification;
  \item it affects the complete recommended signature row without weakening
  a parameter;
  \item it uses only the serialized public key and makes zero signing
  queries;
  \item its output message is fresh and lies in the public map's image; and
  \item it is a design-level failure of the specified raw-vector signature
  interface, rather than a parser, randomness, or implementation defect.
\end{itemize}

\paragraph{Technique attribution.}
The standard attack-goal hierarchy and adaptive chosen-message security
definition allow an adversary to output any fresh message/signature pair,
including a message selected as a function of the public key and the
candidate signature~\cite{hac,gmr}.  We make no novelty claim for this
definition, for choosing a signature before its message, or for public-map
evaluation.  We claim only the concrete attack against TOT's specified
signature relation.

\section{Exact target and claim boundary}

\subsection{Pinned target}

The target is Wuqiang Shen and Shaohua Tang's \emph{TOT, a Fast
Multivariate Public Key Cryptosystem with Basic Secure Trapdoor}, IACR ePrint
2013/771~\cite{tot}.  The current official record has one publicly visible
object, published and last updated on 25 November 2013.  Its 17-page PDF has
SHA-256 digest
\begin{center}
\small\ttfamily
8ee1200e43b60123cf02f0261bee8e381b1229fde8c7722df086d95e6b7784bc.
\end{center}
The extracted text has SHA-256 digest
\begin{center}
\small\ttfamily
4e95403d8f73be474ae0d57ac730370649d4271b2c9dc5b84bd9bb75b4bb794b.
\end{center}

The ePrint gives the signature family $q=2$, $n=128$, $d\geq6$,
$t\geq64$ and reports more than $2^{90}$ work for the best attack considered
there.  The same inventors' signature patent instantiates the boundary row
exactly as
\[
  (q,n,d,t)=(2,128,6,64),
\]
again with a $2^{90}$ security claim~\cite{patent}.  That complete row is the
experimental target.  The attack itself depends only on the public map and
therefore applies independently of the hidden root and affine-map choices.

\subsection{Raw-vector message interface}

Section 3.3 of the ePrint calls
$M=(y'_1,\ldots,y'_n)$ the message to be signed, passes that vector directly
to the trapdoor inversion, and defines verification only by comparison with
the public-map output.  It specifies no document hash, byte encoding, salt,
domain separator, preimage check, or retry wrapper.  The same-inventors
patent independently describes storing the received message, transmitting
the message/signature pair, evaluating the public map, and comparing the
result with the stored message~\cite{patent}.  Its worked four-bit example
passes the vectors directly.

Our claim is therefore about the stated raw message space $\F_2^{128}$.  It
does not assert that the adversary can target an independently chosen
external document or invert a hypothetical document-to-vector hash.  Adding
such a transform would define an additional construction that is absent from
the audited target and would need separate analysis.

No author implementation or normative byte-level public-key, message, or
signature format was located.  The experiment encoding transports the exact
finite-field vectors and complete quadratic map between isolated processes;
it is not claimed as an upstream wire format.

\subsection{Dated attack-history result}

No prior public attack specifically against Shen and Tang's TOT (IACR ePrint
2013/771, sole publicly exposed revision dated 25 November 2013) was located
in searches performed through 13 October 2026.  The search covered the exact
title, ePrint identifier, authors, scheme name, patent family, citations,
repositories, and combinations with attack, cryptanalysis, forgery, and
break terminology.  This is a bounded, dated search result rather than an
absolute priority claim.  Related attacks on multivariate and HFE-family
systems do not disclose this exact-target raw-vector forgery.

\section{The TOT signature relation}

Let $E=\F_{q^n}$ and let $\phi:E\rightarrow\F_q^n$ be the coordinate
isomorphism for a fixed basis.  Key generation chooses distinct
$h_1,\ldots,h_d\in E$ and sets
\[
  g(X)=\prod_{i=1}^{d}(X-h_i),\qquad G(X)=g(X)^t,
\]
where $\gcd(t,q^n-1)=1$.  It also chooses secret invertible affine maps
$S,T:\F_q^n\rightarrow\F_q^n$.  The public map is
\[
  P=T\circ\phi\circ G\circ\phi^{-1}\circ S.
\]
The public key contains the coordinate polynomials of $P$; the trapdoor
contains the roots and affine maps.

Let $v=t^{-1}\bmod(q^n-1)$.  To sign a raw vector $M$, the signer computes
\[
  Y=\phi^{-1}(T^{-1}(M)),\qquad Y_0=Y^v,
\]
finds a root $X$ of $g(X)-Y_0$, and returns
$\sigma=S^{-1}(\phi(X))$.  The verifier accepts exactly when
\begin{equation}
  \Verify(P,M,\sigma)=1 \quad\Longleftrightarrow\quad P(\sigma)=M.
  \label{eq:verify}
\end{equation}

For the experimental row, $t=64=2^6$ and
$v=2^{122}$ because $64\cdot2^{122}=2^{128}\equiv1\pmod{2^{128}-1}$.
In characteristic two,
\[
  g(X)^{64}=\sum_{i=0}^{6}a_i^{64}X^{64i}.
\]
The 64th-power Frobenius is $\F_2$-linear, and the exponents $0,\ldots,6$
have binary weight at most two.  Consequently the central map, and hence
$P$, has Boolean algebraic degree at most two.  Its full constant, linear,
and quadratic coordinate representation occupies
\[
  128\left(1+128+\frac{128\cdot129}{2}\right)
  =1{,}073{,}280
\]
coefficient bits, matching the advertised public-key count.

\section{Zero-query forgery}

\begin{theorem}
For every TOT public key $P:\F_2^{128}\rightarrow\F_2^{128}$ and every
$\sigma\in\F_2^{128}$, the pair $(P(\sigma),\sigma)$ is accepted by the
ordinary verifier.  If the adversary makes no signing query, this pair is a
fresh-message existential forgery against the specified raw-vector
interface.
\end{theorem}

\begin{proof}
The adversary evaluates the public key and defines $M=P(\sigma)$.  Substituting
the pair into Equation~\eqref{eq:verify} gives
$P(\sigma)=M$ identically, so verification accepts.  Since there were no
signing queries, $M$ was not previously submitted to the signing oracle and
is fresh.  For signability, let $u=\phi^{-1}(S(\sigma))$.  After applying
$T^{-1}$ and exponent $v$, the signer obtains
$G(u)^v=g(u)^{tv}=g(u)$, so $u$ belongs to the signer's root set and maps back
to $\sigma$.  Thus $M$ is an honestly signable raw vector, rather than a
value outside the signer's image.
\end{proof}

The experiment fixes a nonzero $\sigma$ before generating the fresh key and
checks that the resulting $M$ is nonzero.  These choices are stronger than
needed for the theorem.  The success probability of the algebraic attack is
one for every key and chosen signature; the nonzero-message check only keeps
the retained transcript easy to inspect.

\section{Full-parameter reproduction}

\subsection{Primary implementation}

The standard-library Python implementation constructs $\F_{2^{128}}$ in a
polynomial basis, samples six distinct roots and two invertible affine maps,
and realizes the exact central and public maps.  It interpolates the complete
quadratic public key into 8,385 coefficient blocks of 128 bits.  Before the
attack, it checks affine inverses, the Frobenius normalization, the signing
equation, and 515 public-versus-secret map evaluations.  It also directly
replays the patent's four-bit worked verifier vector.

Key generation, the public-only attacker, and verification execute as
separate processes.  The generator serializes the public map but no root,
affine map, inverse, random seed, or honest signature.  The attacker receives
only that public map and a sealed empty signing-query transcript.  It chooses
the fixed nonzero $\sigma$, computes $M=P(\sigma)$, and writes the pair.  A
separate ordinary verifier applies Equation~\eqref{eq:verify}.

The coordinator's retained full-row run passed on its first unselected key.
The complete run took 7.93 seconds, used 19,456 KiB peak resident memory and
one worker, and the public evaluation including key parsing took 5.360155
milliseconds.  The ordinary verifier accepted the forged pair.  Four
controls---changed message, changed signature, swapped fields, and a mutated
public key---all rejected.

\subsection{Independent reproduction}

A fresh-context AI agent wrote a separate implementation from the public
specification without inspecting or reusing the primary implementation.  It
generated a new conforming full-row key, gave its forger only public inputs,
used no signing query or secret value, invoked the ordinary verifier, and
repeated independent negative controls.  Its sole sealed full-row run
accepted the forged raw-vector pair and rejected all four controls.  It took
6.798212471 seconds in total, used 20,480 KiB peak resident memory, and used
one worker.  The separately measured public-only attack and ordinary
verification stages each completed in 0.09 seconds.  The public package
contains both implementations and stable wrappers for their complete runs.

\subsection{What the experiment establishes}

The experiment is not a statistical estimate of a rare weakness.  Public
evaluation proves the forgery for every correctly encoded public map; the
full-row executions check the target reconstruction, public/private process
boundary, parameter handling, and verifier behavior.  The fresh keys are not
selected for favorable structure, and no retry is used.

\section{Impact, limitations, and repair direction}

Under the literal interface, TOT does not meet existential unforgeability:
an adversary obtains a fresh accepted pair with one public-map evaluation and
no oracle query.  The attack does not recover the secret key and does not
enable signatures on an independently prescribed raw vector.  Neither fact
is required for existential forgery.

The result does not extend by itself to a construction that first maps an
external document through a preimage-resistant encoding into
$\F_2^{128}$.  Such a wrapper would require the adversary to find a document
whose encoding equals the backward-constructed $P(\sigma)$.  An explicit,
domain-separated, preimage-resistant document-to-vector transform is
therefore a natural direction for blocking this attack.  It would be a
modified scheme and would need its own correctness and unforgeability proof;
we do not claim a complete repair analysis here.

\section{Reproducibility}

The public artifact contains two standard-library Python implementations,
the expected normalized outputs, source pins, and build instructions.  From
the attack directory, run
\begin{verbatim}
python3 reproduce.py
python3 reproduce-all-rows.py
\end{verbatim}
Each command creates a fresh exact-row key in a temporary directory, performs
the public-only attack, invokes the ordinary verifier and controls, removes
ephemeral artifacts, and emits stable JSON.  The package requires one worker
and no network access.

\section*{Acknowledgements}
The computations were performed on the Norwegian Research and Education Cloud (NREC), using resources provided by the University of Bergen and the University of Oslo.

\begin{thebibliography}{9}

\bibitem{tot}
W.~Shen and S.~Tang.
\newblock TOT, a fast multivariate public key cryptosystem with basic secure
trapdoor.
\newblock IACR Cryptology ePrint Archive, Report 2013/771, 2013.
\newblock \url{https://eprint.iacr.org/2013/771}.

\bibitem{patent}
W.~Shen and S.~Tang.
\newblock Multivariate public key signature/verification system and
signature/verification method.
\newblock US Patent 9,948,463 B2, 2018; application publication
US~2016/0234021~A1.
\newblock \url{https://patents.google.com/patent/US9948463B2/en}.

\bibitem{gmr}
S.~Goldwasser, S.~Micali, and R.~L. Rivest.
\newblock A digital signature scheme secure against adaptive chosen-message
attacks.
\newblock \emph{SIAM Journal on Computing}, 17(2):281--308, 1988.

\bibitem{hac}
A.~J. Menezes, P.~C. van Oorschot, and S.~A. Vanstone.
\newblock \emph{Handbook of Applied Cryptography}, Section 11.2.4.
\newblock CRC Press, 1996.
\newblock \url{https://cacr.uwaterloo.ca/hac/}.

\end{thebibliography}

\end{document}
